\section{Canonical forms and all congruences}\label{app:congruences}
\subsection{Canonical forms for several variables}

The small quotients already give compositional function tables at all
arities. There is a second, more syntactic normalization.

Fix a status assignment
$\sigma\in(I\cup\{*\})^n$, where status $i$ means the input is exactly
$\{i\}$, and $*$ means it is ambiguous. On this status cell, represent an
expression by $(C,J)$, denoting
\[
 C\ \cup\ \bigcup_{j\in J}S_j,
 \qquad J\subseteq\{j:\sigma_j=*\}.
\]
For a constant use $(\{i\},\varnothing)$; for a singleton-status variable
use that singleton; for an ambiguous variable use
$(\varnothing,\{j\})$. An owned node collapses to
$(\{r\},\varnothing)$ if either child has that form. Otherwise union the
two $C$ fields and the two $J$ fields. Finally, when $C=I$, discard $J$.

\begin{theorem}[Canonical status forms]\label{view:thm:normal}
This normalization is sound on each status cell. Its final pair $(C,J)$
is uniquely determined by the term function on that cell. Consequently
the table of these pairs over all $(K+1)^n$ statuses is a canonical
representation of any $n$-variable term function.
\end{theorem}
\begin{proof}
If $J$ is nonempty, the value contains an ambiguous input, so cannot be
an owner singleton. If $J$ is empty, the value is exactly $C$. The node
rule therefore implements precisely the two cases of $\star_r$.

For uniqueness, if two $C$ fields differ, choose $k$ in their symmetric
difference and give every ambiguous variable a pair omitting $k$.
Output membership of $k$ then detects the different fields. If $C$
agrees and differs from $I$, but $J$ differs, choose $k\notin C$, give one
variable in the symmetric difference a pair containing $k$, and give all
other ambiguous variables pairs omitting $k$. This separates the
outputs. The only redundant case is $C=I$, already canonicalized by
discarding $J$. All status cells are inhabited for $K\geq3$.
\end{proof}

These rules give a sound and complete normalization decision procedure for term
equality: compute both tables and compare them. They also explain how
ambiguity survives. Once the singleton cases at a cell are fixed, the
remaining computation is just union with a record of which input
variables contribute. The table is not a characterization of which
arbitrary tables are definable by terms, and hence does not yet classify
the full clone. Nor is this a finite unconditional equational basis.

\texttt{structure.py} implements both representations and a complete
counterexample search over the optimal input domain. The status-form
soundness and uniqueness arguments above are not yet mechanized.

\subsection{All congruences and the small irreducible quotients}

Union is term definable, so every algebra congruence is also a union
congruence. The standard correspondence with closure operators is useful
here; the class-maximum construction is classical, compare the lemma to
Proposition 6 in Crapo\cite{Crapo1967}. We give the finite argument
explicitly because the extra owner operations impose a strong condition.

\begin{theorem}[Congruence classification]\label{view:thm:congruences}
The congruences of $\A$ are exactly the kernels of closure operators
$c:\Pow(I)\to\Pow(I)$ satisfying
\begin{equation}\label{view:eq:atomclosure}
 c(\{i\})\in\{\{i\},I\}\qquad(i\in I).
\end{equation}
Equivalently, they correspond to families $\mathcal F\subseteq\Pow(I)$
such that $I\in\mathcal F$, $\mathcal F$ is closed under nonempty
intersections when the result is nonempty, and
\[
 S\in\mathcal F,\ S\ne I,\ i\in S\quad\Longrightarrow\quad
 \{i\}\in\mathcal F.
\]
The closure is the intersection of all members containing the input.
\end{theorem}
\begin{proof}
For a congruence $\theta$, define
$c(S)=\bigcup[S]_\theta$. Finiteness and union compatibility make this
maximum a member of its class. Extensivity and idempotence follow.
If $S\subseteq T$, union compatibility gives
$T\mathrel\theta c(S)\cup T$, so $c(S)\subseteq c(T)$. Thus $c$ is a
closure operator and $\theta=\ker c$.

If $\{i\}$ is identified with any different value $S$, apply
$x\mapsto x\ow i I$. It sends $\{i\}$ to $\{i\}$ and every other input
to $I$. Hence $\{i\}\mathrel\theta I$ and $c(\{i\})=I$.
This crucial step is \texttt{atom\_collapse\_to\_top} in Lean.

Conversely, any closure operator has a union-compatible kernel, since
$c(S\cup T)=c(c(S)\cup c(T))$. Assume~\eqref{view:eq:atomclosure}. Compare
equivalent input pairs for $\star_i$. If neither side absorbs, their
outputs are unions and have equal closures. If both absorb, both outputs
are $\{i\}$. If just one absorbs, an input equivalent to $\{i\}$ differs
from it. Their common closure cannot be $\{i\}$, since extensivity would
make that nonempty input equal to $\{i\}$. Their closure is therefore
$I$, and the nonabsorbing side's union also closes to $I$. This proves compatibility
with every owned operation.

The usual fixed-set description of closure operators gives the final
formulation. A proper closed set containing $i$ rules out
$c(\{i\})=I$ and therefore forces $\{i\}$ to be closed. Conversely, that
condition makes every nonclosed singleton close to $I$.
\end{proof}

In particular, there are exactly $18$ congruences at $K=3$ and $626$ at
$K=4$. The script independently enumerates all partitions at $K=3$.
At $K=4$ it enumerates the classified closure families and checks every
constructed kernel against every operation; completeness uses the proof
above, rather than a partition search.

\begin{theorem}[Subdirectly irreducible homomorphic images]\label{view:thm:si}
Every nontrivial subdirectly irreducible homomorphic image of $\A$ is
obtained from a nonempty proper set $S\subset I$ with closed family
\[
 \mathcal F_S=\{I,S\}\cup\{\{i\}:i\in S\}.
\]
For $|S|=1$ the quotient has two elements. For $|S|\geq2$ it has
$|S|+2$ elements: one singleton state per $i\in S$, the state $S$, and
the state $I$. In particular, every such image has at most $K+1$ elements.
The case $S=I\setminus\{k\}$ is exactly $C_k$.
\end{theorem}
\begin{proof}
The congruence/fixed-family correspondence reverses inclusion.
Meet-irreducible congruences therefore correspond to join-irreducible
admissible fixed families. Every admissible family is the join of the
families $\mathcal F_S$ for its proper members $S$. Each
$\mathcal F_S$ is join-irreducible: its member $S$ cannot be generated by
intersections of its strictly smaller members, which are singletons.
Thus these are exactly the meet-irreducible kernels and hence exactly the
nontrivial subdirectly irreducible quotients.

The displayed fixed sets are representatives of all quotient classes.
For $|S|\geq2$, its least nontrivial congruence identifies the states $S$
and $I$. To see minimality directly, collapsing an atom with $S$ forces
it to collapse with $I$ by its owned operation. Collapsing an atom with
$I$ then collapses $S$ with $I$ by combining it with a different atom
under an owner outside $S$. Collapsing two atoms first collapses one
with $I$ by its owned operation. Every nontrivial collapse therefore
contains the stated pair.
\end{proof}

\begin{corollary}[Minimum subdirect decomposition]
Every subdirect representation of $\A$ by nontrivial subdirectly
irreducible factors must include, up to the quotient maps' kernel,
each $C_k$. The $K$-factor decomposition of
Corollary~\ref{view:cor:product} is consequently minimum and irredundant.
\end{corollary}
\begin{proof}
The distinct values $I$ and $I\setminus\{k\}$ must be separated by some
factor. Among the classified quotients, only
$\mathcal F_{I\setminus\{k\}}$ separates them: every other proper
fixed set fails to contain $I\setminus\{k\}$, so that value closes to
$I$. Thus the corresponding factor $C_k$ is necessary for every $k$.
\end{proof}

These statements concern homomorphic images of $\A$ itself. They do not
bound subdirectly irreducible members of the whole variety generated by
$\A$, and do not establish finite basability.

\subsection{What internal role switching changes}

Now expand the signature by unary operations $S\mapsto\pi S$ for every
role permutation $\pi$. Write $\A^{\mathrm{sym}}$ for this algebra. This
is a separate extension: the preceding small test-domain and one-player
quotient theorems used the original signature.

\begin{theorem}[The symmetric congruence chain]\label{view:thm:sym}
There are exactly $K$ congruences of $\A^{\mathrm{sym}}$. They are the
kernels of
\[
 c_m(S)=\begin{cases}S,&|S|\leq m,\\I,&|S|>m,\end{cases}
 \qquad m=0,1,\ldots,K-1.
\]
Here $m=0$ gives the universal relation and $m=K-1$ gives equality.
For $1\leq m\leq K-1$, the quotient has
$1+\sum_{j=1}^{m}\binom Kj$ values. These are the restrictions to
nonempty sets of the uniform-matroid closures of rank $m+1$.
\end{theorem}
\begin{proof}
A congruence of the expansion is exactly a congruence of $\A$ invariant
under every permutation. Its class-maximum closure therefore commutes
with permutations. If any singleton closes to $I$, symmetry makes every
singleton close to $I$, and monotonicity makes every nonempty set close
to $I$.

Otherwise all singletons are closed. Let $m$ be the largest size of a
proper closed set. Permutation invariance makes all $m$-element sets
closed. Their nonempty intersections give every nonempty set of size at
most $m$: for a smaller set $T$ and $x\notin T$, an $m$-element superset
of $T$ omitting $x$ exists because $m<K$. Thus the proper closed sets
are exactly those of size at most $m$. The congruence classification
gives precisely the displayed kernels. Their closures are nested, so
their kernels form a chain. The cardinalities and the uniform-matroid
identification follow from their fixed sets; see the uniform matroids
in Oxley\cite{Oxley2011}.
\end{proof}

\begin{corollary}[Symmetry makes the algebra irreducible]
$\A^{\mathrm{sym}}$ is subdirectly irreducible. Its least nontrivial
congruence identifies all sets of size at least $K-1$, leaving every
smaller nonempty set fixed.
\end{corollary}
\begin{proof}
In the chain above, $\ker c_{K-2}$ is the least congruence strictly above
equality. It exists for $K\geq3$ and has the stated classes.
\end{proof}

The quotient with $m=1$ records every unique winner separately and merges
all ambiguity. The intervening quotients retain increasingly large sets
of alternative equilibrium winners. Thus internal role symmetry permits
an exact hierarchy of ambiguity thresholds.

The one-player contextual quotient also changes: it is now equality.
If $S,T$ differ at role $i$, choose a permutation taking $i$ to $k$.
The context $\pi(x)$ then distinguishes their possible-win observation
for $k$. This does not invalidate the earlier product representation:
the tuple of all $C_k$ still represents the underlying algebra, and
permutations act by transporting \emph{between} its coordinates. They
cannot be interpreted as separate unary operations inside each fixed
factor. The distinction between external symmetry and internal role
switching therefore has a measurable algebraic effect.
